\section{Anatomy of Algorithmic Failures}
\label{sec:cases}

To ground the economic model in Section~\ref{sec:tco}, we analyze three
landmark cases of algorithmic failure between 2018 and 2022. We examine
these not as ethical lapses in isolation, but as \emph{structural
failures of accountability mechanisms}: the inability to detect, persist,
or audit decisions at the point of impact. Each failure maps to a
specific gap addressed by the BTV stack (Papers~1--3~\cite{btv2026,btv-t-repo,btv-ar2026}).

\subsection{Case 1: Racial Bias in Healthcare Allocation}
\label{subsec:obermeyer}

\noindent\textit{System:} Optum Impactability Algorithm (US health system, $\sim$200M patients).\\
\textit{Decision:} Allocation of high-risk care management resources.\\
\textit{Failure mode:} Proxy discrimination via silent evidence.

Obermeyer et al.~\cite{obermeyer2019dissecting} revealed that the
algorithm systematically privileged white patients over Black patients
because it used \emph{healthcare costs} as a proxy for \emph{health
need}. Since less money is historically spent on Black patients for the
same severity of illness, the proxy produced a structurally biased
rankings. The bias was invisible at the decision layer because the
proxy was never validated against the regulatory intent it was meant to
operationalise.

\textbf{BTV gap --- Paper~1 (Evidence Binding):} The decision logic was
not bound to an evidence token that verified the semantic validity of the
input proxy. Under the Paper~1 invariant, a
$\texttt{Verdict} \multimap E \otimes C$ type constraint would require
the input features to be validated against a declared compliance policy
before a verdict can be constructed. A cost-based feature used as a
health-need proxy would fail this validation at compile time, surfacing
the bias before deployment rather than two years into production.

\textbf{Cost:} Optum revised the algorithm following publication;
latent litigation and reputational exposure remained unquantified at time
of writing. The case produced no direct regulatory fine, illustrating
that even ``self-corrected'' opacity carries cost in the form of
prolonged harm duration.

\subsection{Case 2: Gender Bias in Automated Hiring}
\label{subsec:amazon}

\noindent\textit{System:} Amazon internal resume screening tool (2014--2018).\\
\textit{Decision:} Candidate ranking and rejection.\\
\textit{Failure mode:} Silent decisions; opaque rejections.

Reuters reported in 2018 that Amazon scrapped an AI recruiting tool
that penalised resumes containing the word ``women's'' or originating
from women's colleges~\cite{dastin2018amazon}. The model had been trained
on a decade of resumes submitted predominantly by men, and learned
gender as a latent penalisation signal. Rejected candidates received no
auditable rationale.

\textbf{BTV gap --- Paper~1 (Evidence Binding):} The system produced
rejections without binding each decision to a specific evidence token.
Under the Paper~1 invariant, a rejection verdict cannot be constructed
without an evidence token that specifies the operative features. The
bias would have been visible to auditors at the first audit cycle:
rejection receipts citing gender-correlated features would produce a
statistically anomalous evidence distribution.

\textbf{BTV gap --- Paper~2 (Persistence):} Amazon terminated the
project in 2018, effectively deleting the operational record of the
bias mechanism. A Paper~2 append-only Transparency Log would have
preserved the full decision history for regulatory review regardless of
project lifecycle. Voluntary cessation of an opaque system is not a
substitute for verifiable accountability.

\textbf{Cost:} Amazon wrote off the R\&D investment and avoided
regulatory fines by ceasing operations before formal investigation. This
represents the \emph{opacity-as-retreat} strategy: opacity creates the
option to quietly discontinue a discriminatory system before legal
exposure crystallises. The BTV stack closes this option structurally ---
the record exists regardless of whether the system continues to operate.

\subsection{Case 3: Housing Discrimination via Ad Delivery}
\label{subsec:meta}

\noindent\textit{System:} Meta ``Special Ad Audiences'' algorithm.\\
\textit{Decision:} Targeting of housing, credit, and employment advertisements.\\
\textit{Failure mode:} Selective delivery; black-box persistence.

In 2022, the US Department of Justice settled with Meta over allegations
that its ad delivery algorithms violated the Fair Housing Act by
excluding protected demographic classes even when advertisers specified
neutral targeting~\cite{doj2022meta}. The algorithm optimised for
engagement in ways that produced discriminatory delivery patterns;
because decisions occurred at millions of impressions per day, no
external auditor could inspect individual routing decisions.

\textbf{BTV gap --- Paper~2 (Persistence):} Meta maintained internal
logs, but the causal reasoning for why a specific user received or was
excluded from an impression was not externally accessible. Regulators
were forced to rely on statistical disparity testing rather than direct
causal reconstruction --- a weaker evidentiary standard that took four
years of litigation to resolve. A public \texttt{InclusionReceipt} per
impression batch would have reduced investigation time to the duration
of a cryptographic verification.

\textbf{BTV gap --- Paper~3 (Privacy-Preserving Audit):} Meta could not
publish targeting data without violating the privacy of individual users.
This created a genuine tension between transparency and data protection
that the DOJ settlement did not resolve technically. A Paper~3
zero-knowledge consistency proof would have allowed Meta to demonstrate
statistical non-discrimination across protected attributes without
revealing any individual's viewing or demographic profile.

\textbf{Cost:} Meta paid no civil monetary penalty in the 2022
settlement but accepted drastic restrictions on ad-targeting capabilities
--- a recurring revenue sacrifice estimated at \$200M/year~\cite{meta-settlement-analysis} --- and committed to ongoing third-party audits.
The absence of a verifiable audit layer converted a tractable compliance
problem into a multi-year enforcement action with permanent product
constraints.

\subsection{Synthesis: Opacity as a Structural Choice}
\label{subsec:synthesis}

Table~\ref{tab:case-gaps} maps the three cases to the BTV stack layers
they would have engaged.

\begin{table}[h]
\centering
\caption{Case Studies Mapped to BTV Accountability Gaps}
\label{tab:case-gaps}
\begin{tabular}{lccc}
\toprule
\textbf{Case} & \textbf{Paper~1} & \textbf{Paper~2} & \textbf{Paper~3} \\
              & \footnotesize{Evidence binding} & \footnotesize{Persistence} & \footnotesize{ZK audit} \\
\midrule
Obermeyer (Optum, 2019)  & \checkmark & --- & --- \\
Amazon hiring (2018)     & \checkmark & \checkmark & --- \\
Meta ad delivery (2022)  & --- & \checkmark & \checkmark \\
\bottomrule
\end{tabular}
\end{table}

Across all three cases, opacity was not an accident but a structural
property of the deployment architecture: no layer required the system to
produce verifiable, persistent, auditable evidence of its own behaviour.
The economic question --- whether prevention would have been cheaper than
remediation --- is addressed quantitatively in
Section~\ref{sec:forensic}, where we compute the Cost of Opacity ratio
for 20 enforcement decisions. We show that in every case, the fine or
sanction paid exceeded the annual cost of the BTV infrastructure by
between 18$\times$ and 40{,}000$\times$.
